\chapter{Stock Markets and Equity Claims}
\label{chap:Stock_Markets_and_Equity_Claims}

\begin{openingcase}{Why Can Market Value Change When No Cash Changes Hands?}

A listed company has 100 million shares outstanding, and its shares trade at RMB~10. The company's market value of equity is RMB~1 billion. One morning, the company announces that a new product has passed an important technical milestone. No dividend is paid, no new shares are issued, and no cash enters the company. Yet investors revise their expectations about future profits, and the share price rises to RMB~12.

The company's market value of equity has increased by RMB~200 million even though the company did not receive RMB~200 million of new financing. The change occurred in the secondary market, where investors exchanged existing shares with one another.

\end{openingcase}

\begin{tcolorbox}[
  enhanced,
  colback=FinanceBlue!6,
  colframe=FinanceBlue,
  boxrule=0.6pt,
  sharp corners,
  left=8pt,
  right=8pt,
  top=7pt,
  bottom=7pt,
  before skip=1em,
  after skip=1.6em,
  title=\textbf{Core Questions},
  coltitle=white,
  colbacktitle=FinanceBlue,
  fonttitle=\bfseries,
  attach boxed title to top left={xshift=6pt,yshift=-2mm},
  boxed title style={
    boxrule=0pt,
    sharp corners,
    left=5pt,
    right=5pt,
    top=2pt,
    bottom=2pt
  }
]
\begin{enumerate}[
  leftmargin=*,
  label=\textcolor{FinanceBlue}{\textbf{\arabic*.}},
  itemsep=0.55em,
  topsep=0.2em
]

\item What rights does an equity claim provide, and
how do residual cash-flow rights, limited liability,
and control differ from the rights attached to debt?

\item How do ownership structure, voting rights, and
corporate governance shape the allocation of control
and agency conflicts among shareholders and managers?

\item How are equity claims issued and traded in stock
markets, and how are information, growth expectations,
and risk reflected in market prices and required returns?

\end{enumerate}
\end{tcolorbox}

\section{Equity as a Residual Ownership Claim}

Equity is a residual claim on a corporation. Shareholders do not hold a contractual promise to receive fixed payments. They hold rights whose value depends on what remains after the firm's contractual obligations have been met and on the governance rights attached to the shares.

\subsection{Residual Cash-Flow Rights and Limited Liability}

A corporation is a separate legal entity. The shares
carry financial and governance rights defined by
corporate law, the firm's constitutional documents,
and the terms of the relevant share class. Common
equity carries a residual claim on corporate cash
flows and assets. In liquidation, common shareholders
rank behind debt holders and other senior claimants.

For fully paid shares, limited liability generally
means that shareholders are not required to contribute
additional personal assets solely by virtue of their
share ownership. Separate obligations arising from
personal guarantees, borrowing, derivatives, or other
contracts remain enforceable.

Residual status does not imply immediate distribution.
A profitable corporation may retain earnings and
reinvest them rather than pay them out. Shareholders
retain a residual claim on the value generated by
those resources, subject to future business outcomes
and the governance arrangements governing their use.
Dividends are one means of distributing value to
shareholders, not the definition of the residual claim.

\subsection{Equity Financing and Risk Bearing}

Equity financing provides funds without creating fixed
contractual interest or principal payments. A company
that issues common shares raises capital in exchange
for claims on future residual cash flows and, usually,
voting rights. This payment flexibility can be valuable
when operating cash flows are uncertain.

The absence of fixed payments is accompanied by
residual risk bearing. If operating performance
deteriorates, common dividends can be reduced or
omitted and the value of equity can fall sharply.
If performance is strong, shareholders participate
in the upside after senior claims are satisfied.
Common equity therefore bears much of the variation
in firm value after contractual claims have been
taken into account.

This risk-bearing structure can support investment
in uncertain projects. The underlying business risk
remains, but the financing structure affects how
adverse operating outcomes are transmitted to
different claimants.

\begin{example}{Business Risk and Risk Bearing}
{business-risk-bearing}

Consider a firm developing a new technology with an
initial investment of RMB~100 million, financed by
RMB~40 million in debt and RMB~60 million in equity.
Suppose the debt contract requires repayment of
RMB~44 million at maturity. Assume that the amounts
generated by the project represent the total value
available to meet these claims, with no other assets
or liabilities.

If the project generates RMB~140 million at maturity,
debt holders receive the contracted RMB~44 million,
and shareholders receive the remaining RMB~96 million.
If project value falls to RMB~60 million, debt holders
still receive RMB~44 million, leaving RMB~16 million
for shareholders. If project value falls further to
RMB~30 million, shareholders lose their entire
investment, and debt holders recover only RMB~30 million.


\end{example}

\subsection{Debt versus Equity: Priority, Control, and Risk Allocation}

Debt and equity are different claims on the same underlying business. Their distinction extends beyond the form of payment. They differ in priority, control rights, and exposure to changes in firm value.

Debt gives creditors contractually specified claims and priority over common equity. For ordinary debt, the upside is generally limited to the contractual payoff, while losses arise when promised payments are not made or recovered in full. Creditors usually do not exercise general corporate control, but covenants, enforcement rights, and restructuring arrangements can give them greater influence when contractual conditions are breached or the firm enters financial distress.



\begin{table}[htbp]
\centering

\caption{Debt and Common Equity as Financial Claims}
\label{tab:ch9-debt-equity-claims}

\begin{tabularx}{\textwidth}{@{}p{3.0cm}Y Y@{}}

\toprule

\textbf{Dimension}
& \textbf{Debt}
& \textbf{Common equity} \\

\midrule

Cash-flow rights
& Contractually specified interest and/or principal payments, subject to default
& Residual claim on corporate cash flows after contractual obligations are met \\

Priority
& Generally ranks ahead of common equity; priority among debt claims depends on seniority, security, and applicable law
& Residual claimant in liquidation \\

Maturity
& Generally specified, although perpetual debt also exists
& No contractual maturity \\

Control
& Generally lacks ordinary voting rights; contractual protections and enforcement rights apply as specified in the debt agreement
& Voting and governance rights, subject to the share class and applicable law \\

Upside
& Generally limited to contractual payments
& Participates in residual gains and increases in corporate value \\

Downside
& May incur losses if contractual claims are not fully recovered
& Investment can be lost entirely, but shareholders are generally not personally liable for corporate debts solely by virtue of share ownership \\

\bottomrule

\end{tabularx}

\end{table}

Common equity is a residual claim. Shareholders have no promised maturity payment and receive the value remaining after senior claims have been taken into account. Subject to limited liability, the value of the share can fall to zero, while the upside is not capped by a contractual repayment amount. Under normal conditions, voting rights also give common shareholders governance rights, including the election of directors and approval of specified corporate matters.

These differences can create conflicts between creditors and shareholders. Once debt is outstanding, shareholders may prefer projects with greater upside even when part of the additional downside is borne by creditors. Creditors respond through pricing, collateral, covenants, maturity provisions, and other contractual protections. In financial distress, bargaining power and control can also shift toward creditors.


\section{Ownership, Control, and Corporate Governance}

The corporate form separates ownership from day-to-day management. Shareholders hold residual ownership claims, while directors and managers exercise delegated control over corporate resources. This separation can create agency problems when decision makers do not bear the full consequences of their decisions. Corporate governance allocates control rights and provides mechanisms for monitoring and disciplining management.

\subsection{Voting Rights and Board Control}

Shareholders generally do not make routine operating decisions. They exercise control primarily through voting rights, including the election of directors and approval of specified major matters. The board appoints and oversees senior management and exercises the decision rights assigned to it under corporate law and the firm's governance arrangements. The board is therefore the principal link between shareholder voting rights and managerial control.

A simplified chain of delegated control is shown in Figure~\ref{fig:ch9-governance-chain}.

\begin{figure}[htbp]
\centering

\begin{tikzpicture}[
  >=Latex,
  box/.style={
    draw,
    rounded corners,
    align=center,
    minimum width=4.2cm,
    minimum height=0.85cm,
    inner xsep=6pt,
    inner ysep=4pt,
    font=\normalsize
  },
  flow/.style={
    -Latex,
    line width=0.85pt,
    FinanceBlue
  },
  lab/.style={
    font=\small,
    text=FinanceBlue,
    fill=white,
    inner xsep=3pt,
    inner ysep=1pt
  },
  node distance=0.75cm
]

\node[box] (share) {Shareholders};

\node[box, below=of share] (board)
  {Board of directors};

\node[box, below=of board] (mgmt)
  {Management};

\node[box, below=of mgmt] (assets)
  {Corporate assets and operations};

\draw[flow]
  (share) --
  node[lab, right] {elect directors}
  (board);

\draw[flow]
  (board) --
  node[lab, right] {appoint and oversee}
  (mgmt);

\draw[flow]
  (mgmt) --
  node[lab, right] {manage}
  (assets);

\end{tikzpicture}

\caption{A simplified chain of delegated control}
\label{fig:ch9-governance-chain}
\end{figure}

Under a one-share, one-vote structure, formal voting
power is proportional to the number of voting shares
held. A shareholder holding 30\% of the
outstanding voting shares therefore holds 30\%
of the associated voting rights.Effective influence, however, also depends on ownership concentration, voter participation, voting agreements, board-election rules, and the matters subject to shareholder approval.

Voting rights have economic value when they enable
shareholders to influence corporate decisions.
A shareholder with a minority economic interest
may still exercise substantial influence when
ownership is widely dispersed and other shareholders
participate less actively in voting. The distribution
of voting rights therefore need not translate directly
into the distribution of effective control.

\subsection{Agency Problems and Ownership Structure}

Agency problems arise when decision authority is delegated and the interests of decision makers and claimants do not fully coincide. In listed corporations, ownership structure affects both the incentives to monitor corporate decisions and the conflicts that may arise from control. Two relationships are especially important.

The first is between managers and shareholders. Managers control corporate resources but may hold only a small fraction of the residual cash-flow rights. Their objectives may therefore differ from those of shareholders. Managers may value private benefits, organizational expansion, greater job security, or lower personal risk even when these choices reduce shareholder value.

Dispersed ownership can weaken monitoring incentives. A small shareholder may bear much of the cost of investigating management while receiving only a small fraction of any resulting improvement in firm value. Other shareholders can benefit from the same monitoring without bearing its cost. This free-rider problem can reduce individual incentives to monitor and increase the importance of boards, institutional investors, disclosure, auditing, and other governance mechanisms.

The second relationship is between controlling and
minority shareholders. Concentrated ownership changes
monitoring incentives. A large shareholder captures
a greater share of the benefits from improved
monitoring and may therefore have a stronger
incentive to oversee management. Greater concentration,
however, can create conflicts between controlling
and minority shareholders. A controlling shareholder
may influence related-party transactions, asset
transfers, financing decisions, or other corporate
actions in ways that provide private benefits while
reducing value for minority investors.


\begin{table}[H]

\centering

\caption{Two important equity agency problems}

\label{tab:ch9-agency-problems}

\begin{tabularx}{\textwidth}{@{}p{3.0cm}Y Y@{}}

\toprule

\textbf{Agency problem} & \textbf{Source of conflict} & \textbf{Potential governance responses} \\

\midrule

Managers versus shareholders
& Managers exercise decision authority but may hold only part of the residual claim
& Board oversight, compensation design, disclosure, monitoring, voting, takeover pressure \\

Controlling versus minority shareholders
& Controlling shareholders can influence decisions whose effects are shared with minority investors
& Related-party transaction rules, independent directors, disclosure, voting safeguards, legal protection \\

\bottomrule

\end{tabularx}

\end{table}


\subsection{Shareholder Rights and Share Classes}

A share represents a bundle of financial and governance
rights. The rights attached to a particular share depend
on corporate law, the firm's governing documents, and
the relevant share class. Differences across share
classes can affect both the distribution of financial
benefits and the allocation of voting power.

Share classes can allocate these rights differently. Preferred shares, for example, may have priority over common shares in dividends or liquidation while carrying limited or no voting rights. They remain junior to creditor claims. Conversion, redemption, or other contractual features can give preferred shares characteristics associated with both equity and debt.

Dual-class and weighted-voting structures separate voting power from economic ownership more directly. Let class \(j\) have \(N_j\) shares outstanding and \(v_j\) votes per share. Its fraction of total voting power is

\begin{equation}
\omega_j
=
\frac{v_jN_j}{\sum_k v_kN_k}.
\label{eq:ch9-voting-power}
\end{equation}

If all classes have the same cash-flow rights per share, class \(j\)'s fraction of cash-flow ownership is

\begin{equation}
\alpha_j
=
\frac{N_j}{\sum_k N_k}.
\label{eq:ch9-cash-flow-share}
\end{equation}

When voting rights differ across classes, \(\omega_j\) and \(\alpha_j\) need not coincide.

Such structures can allow founders or other controlling
shareholders to retain voting control while raising
equity from outside investors. This may support
continuity in corporate strategy, particularly when
investment projects have long horizons. At the same
time, the separation can weaken accountability to
outside shareholders whose economic exposure may
exceed their voting influence. The governance
implications depend on the safeguards provided by
corporate law and listing rules.

\begin{tcolorbox}[
  enhanced,
  breakable,
  colback=white,
  colframe=FinanceBlue,
  boxrule=0.7pt,
  sharp corners,
  borderline west={2.2pt}{0pt}{FinanceBlue},
  title=\textbf{Financial Perspective: Cash-Flow Rights and Voting Rights Need Not Coincide},
  coltitle=white,
  colbacktitle=FinanceBlue,
  fonttitle=\bfseries,
  before skip=1.2em,
  after skip=1.2em
]

A class of shares may account for a relatively small
fraction of a firm's total cash-flow rights while
carrying a larger fraction of its voting power.
Cash-flow rights determine participation in residual
economic value, whereas voting rights determine the
ability to influence corporate decisions.

\end{tcolorbox}

\subsection{Corporate Actions and Changes in Ownership Claims}

Corporate actions can change the cash held by the firm, the number of shares outstanding, or the distribution of rights among shareholders. Their immediate effects can be analyzed by considering whether cash enters or leaves the firm, whether shares are created or retired, and whether ownership fractions or rights change.

A new equity issue increases the number of shares outstanding and can change existing ownership fractions. Suppose a company has \(N_0\) shares outstanding and issues \(N_1\) new shares. An investor who owns \(n\) shares before the issue holds a fraction $\alpha_0=n/N_0$. If the investor does not participate in the new issue, the ownership fraction becomes $\alpha_0=n/(N_0+N_1)$  

The investor's percentage ownership is diluted. Percentage dilution, however, does not necessarily imply a loss of economic value. If new shares are issued at a fair price, the cash received by the firm increases its value by approximately the amount raised at the time of issuance. A smaller ownership fraction can therefore apply to a larger total equity value. If shares are issued at an unusually low price to selected investors, value can instead be transferred from nonparticipating shareholders to the new investors.

A rights issue allows existing shareholders to purchase new shares and can help them preserve their proportional ownership if they participate. Shareholders who do not participate may experience percentage dilution; transferable rights can provide compensation for the value of the subscription privilege.

\begin{example}{Percentage Dilution and Value Transfer}{new-equity-dilution}

A company has 80 million shares outstanding, trading at RMB~10 per share, giving a market equity value of RMB~800 million. An investor owns 8 million shares, or 10\% of the company.

The company issues 20 million new shares, and the investor does not participate. The investor's ownership fraction falls to

\[
\frac{8}{100}=8\%.
\]

If the new shares are issued at RMB~10, the company raises RMB~200 million. Assuming no other change in value, post-issue equity value becomes RMB~1 billion and the share price remains RMB~10. The investor's percentage ownership falls, but the value of the 8 million shares remains RMB~80 million.

If the new shares are instead issued to selected investors at RMB~5, the company raises only RMB~100 million. Under the same assumptions, post-issue equity value is approximately RMB~900 million, implying a share value of

\[
\frac{900}{100}=\text{RMB }9.
\]

The investor's shares are then worth approximately RMB~72 million. In this case, the below-market issue transfers value from existing shareholders to the new investors.

Percentage dilution and loss of economic value are therefore distinct. The economic effect depends on both the change in ownership fractions and the issue price.

\end{example}

Dividends and share repurchases both distribute cash from the firm to shareholders, but they affect ownership claims differently. A cash dividend transfers resources from the corporation to shareholders without changing the number of shares outstanding. If a firm pays a dividend \(D\) per share and no other relevant information arrives, the share price should fall when the share begins trading without the right to that dividend. Abstracting from taxes and market frictions,

\begin{equation}
P_{\mathrm{ex}}
\approx
P_{\mathrm{cum}}-D,
\label{eq:ch9-ex-dividend-price}
\end{equation}

where \(P_{\mathrm{cum}}\) is the price with the dividend right and \(P_{\mathrm{ex}}\) is the ex-dividend price. The shareholder receives cash while the firm's cash holdings decline.

A share repurchase also transfers cash from the firm to shareholders, but it reduces the number of shares outstanding. Shareholders who sell receive cash, while shareholders who retain their shares own a larger fraction of the remaining equity, all else equal. A repurchase can affect value through the price paid, information conveyed by the transaction, taxes, or changes in capital structure. A reduction in the number of shares alone does not create economic value.

A stock split changes the number of shares without transferring cash or changing proportional ownership. In a \(m\)-for-\(n\) split, the number of shares becomes

\begin{equation}
N'=\frac{m}{n}N,
\label{eq:ch9-stock-split-shares}
\end{equation}
and, absent other effects, the price per share adjusts approximately to

\begin{equation}
P'=\frac{n}{m}\times P.
\label{eq:ch9-stock-split-price}
\end{equation}

Total market capitalization therefore remains unchanged if the split conveys no new information and does not otherwise affect value.

\begin{table}[H]
\centering

\caption{Selected Corporate Actions and Their Immediate Claim-Level Effects}
\label{tab:ch9-corporate-actions}

\begin{tabularx}{\textwidth}{@{}p{2.8cm}Y Y@{}}

\toprule

\textbf{Action}
& \textbf{Immediate change}
& \textbf{Key economic question} \\

\midrule

Cash dividend
& Cash is distributed to shareholders; shares outstanding remain unchanged. The share price typically adjusts downward on the ex-dividend date.
& How does the distribution affect the resources retained by the firm and shareholder value? \\

Share repurchase
& Cash leaves the firm and shares outstanding generally decline.
& At what price are shares repurchased, and how are ownership fractions and shareholder value affected? \\

Stock split
& Shares outstanding increase and the price per share adjusts proportionally; total equity value and ownership fractions remain unchanged, other things equal.
& Does the split convey information or affect trading conditions? \\

New equity issue
& Cash enters the firm and shares outstanding increase, potentially diluting existing shareholders' ownership percentages.
& At what price are new shares issued, and how does the issue affect existing shareholders' economic value? \\

\bottomrule

\end{tabularx}

\end{table}

\section{Stock Market Architecture and Price Discovery}

Stock markets perform two linked functions. Primary markets allow firms to raise equity capital, while secondary markets allow existing shares to be transferred among investors. In secondary trading, competing buy and sell orders based on investors' information and expectations generate market prices, supported by exchanges, brokers, clearing systems, and depositories.

\subsection{Equity Issuance and Primary Markets}

A company raises new equity capital in the primary market by issuing shares to investors. An initial public offering (IPO) offers shares to public investors for the first time, while a seasoned equity offering raises additional public equity after listing. Rights issues and private placements provide other channels for issuing new shares.

Equity issuance requires the terms and price of the new claim to be determined. The issuer decides how many shares to offer and what rights they carry, while investors assess the firm's expected cash flows, risks, governance, and use of proceeds. Intermediaries such as underwriters, bookrunners, and financial advisers can support due diligence, disclosure, bookbuilding, and placement. Their role is to facilitate the issuance process and match new shares with investor demand.

A public offering may contain both primary and secondary shares. Primary shares are newly issued by the company and bring new funds into the firm. Secondary shares are sold by existing shareholders and transfer ownership without providing sale proceeds to the company. The distinction is important when interpreting the amount of capital raised by an offering.

\begin{figure}[H]
\centering

\begin{tikzpicture}[
  >=Latex,
  box/.style={
    draw,
    rounded corners=2pt,
    line width=0.7pt,
    align=center,
    text width=1.9cm,
    minimum height=0.9cm,
    inner sep=5pt
  },
  cash/.style={
    -Latex,
    line width=0.9pt,
    FinanceBlue
  },
  claim/.style={
    -Latex,
    dashed,
    line width=0.9pt,
    FinanceBlue
  },
  lab/.style={
    font=\small,
    align=center,
    fill=white,
    inner sep=1.5pt
  }
]

% Participants
\node[box] (firm) at (-1.0,0)
  {Company};

\node[box] (market) at (4.5,0)
  {Primary market};

\node[box] (investors) at (10.0,0)
  {Investors};

% New equity claims: company to investors
\draw[claim]
  ([yshift=5pt]firm.east) --
  node[lab, above=3pt] {New equity claims}
  ([yshift=5pt]market.west);

\draw[claim]
  ([yshift=5pt]market.east) --
  node[lab, above=3pt] {New equity claims}
  ([yshift=5pt]investors.west);

% Capital: investors to company
\draw[cash]
  ([yshift=-5pt]investors.west) --
  node[lab, below=3pt] {Capital}
  ([yshift=-5pt]market.east);

\draw[cash]
  ([yshift=-5pt]market.west) --
  node[lab, below=3pt] {Capital}
  ([yshift=-5pt]firm.east);

\end{tikzpicture}

\caption{Capital and Equity Claim Flows in the Primary Market}
\label{fig:ch9-primary-equity}
\end{figure}

The offer price affects the amount of capital raised
and the allocation of ownership between existing and
new shareholders. If the price is set too high relative
to investors' valuations, the offering may attract
insufficient demand. If it is set too low, the issuer
may raise less capital than it could otherwise obtain.
An initial pricing discount may reflect valuation
uncertainty, information asymmetry, and placement risk.

\subsection{Exchanges and Alternative Trading Venues}

An exchange provides rules and infrastructure for listing and trading securities. Listing and trading are related but distinct. A company may be listed on one exchange while transactions in its shares also occur through other permitted trading venues.

Organized exchanges bring orders together under
common trading rules, access arrangements, and
market-data systems. Modern equity markets may
also include alternative trading systems, such as
crossing networks and dark pools, as well as
dealer-based trading venues.

Consider a company whose shares are listed on the
New York Stock Exchange (NYSE). Investors do not
necessarily have to trade these shares on the NYSE.
Depending on applicable rules and order-routing
arrangements, a broker may execute an order on
another exchange or through an alternative trading
system (ATS), such as a dark pool. The same shares
can therefore be traded across multiple venues,
each with its own order flow and available liquidity.

Multiple venues can increase competition in execution services and allow different trading mechanisms to coexist. They can also fragment order flow, making available liquidity more difficult to observe across the market. Venue structure therefore affects how orders interact, where liquidity is concentrated, and how prices are formed.

Trading venues operate within a broader market
infrastructure. Brokers route orders, clearing
arrangements establish and, where applicable,
net the obligations arising from trades, while
securities depositories maintain ownership records
and settlement systems complete the transfer of
cash and securities. Electronic trading has changed
the speed and organization of these processes
without changing their basic economic functions.

\subsection{Trading Mechanisms, Orders, and Market Liquidity}

A trading mechanism determines how buy and sell orders interact. In an order-driven market, orders are submitted to a limit order book and transactions occur when compatible orders meet. In a dealer market, dealers quote prices at which they are willing to buy and sell. Many modern equity markets combine elements of these structures.

A market order gives priority to execution at the best available prices, subject to available quantity. A limit order specifies the highest price a buyer is willing to pay or the lowest price a seller is willing to accept. It provides greater price control but may remain unexecuted. Market orders typically demand available liquidity, while nonmarketable limit orders can supply liquidity to other traders.

If the best bid and ask are \(P_{\mathrm{bid}}\) and \(P_{\mathrm{ask}}\), the quoted bid--ask spread is

\begin{equation*}
s=P_{\mathrm{ask}}-P_{\mathrm{bid}}.
\label{eq:ch9-bid-ask-spread}
\end{equation*}

The midpoint is

\begin{equation*}
M=\frac{P_{\mathrm{ask}}+P_{\mathrm{bid}}}{2},
\label{eq:ch9-midquote}
\end{equation*}
and the proportional quoted spread is

\begin{equation*}
s_{\mathrm{rel}}
=
\frac{P_{\mathrm{ask}}-P_{\mathrm{bid}}}{M}.
\label{eq:ch9-relative-spread}
\end{equation*}

The spread is one dimension of liquidity. Depth measures the quantity available near current prices, while price impact describes how strongly execution prices respond to order size. A market can display a narrow spread for small trades while offering limited depth for larger orders.

The supply of liquidity can vary with market conditions. When uncertainty or adverse-selection risk increases, liquidity suppliers may quote wider spreads, display less depth, or withdraw orders. The possibility of trading against better-informed investors therefore links information conditions to market liquidity.

Trading constraints also affect how information and investor views enter the order flow. Short selling allows an investor to borrow shares, sell them, and later repurchase shares to return to the lender. It enables negative information or expectations to be expressed through sell orders even when the investor does not initially own the shares. Restrictions or high borrowing costs can limit this channel. Short selling can therefore affect order flow, liquidity, and price discovery, although its effects depend on market conditions and trading rules.

\subsection{Information, Prices, and Price Discovery}

A stock price is an observed trading price for an equity claim. It can change when investors revise expectations about future cash flows or required returns, or when the rights attached to the claim change.

Let \(\mathcal{I}_t\) denote the information available at time \(t\). A dividend-based valuation benchmark can be written as

\begin{equation}
P_t
=
\sum_{j=1}^{\infty}
\frac{\mathrm{E}[D_{t+j}\mid\mathcal{I}_t]}
{(1+r_{e,t})^j},
\label{eq:ch9-information-price}
\end{equation}
where \(r_{e,t}\) is treated as a constant required return from time \(t\) onward. New information can change expected future dividends, the required return, or both, and therefore change the value investors assign to the share.

Price discovery is the process through which information and investor assessments are reflected in market prices through trading. Investors may differ in their information, models, expectations, and liquidity needs, so their orders reveal only part of what they know. Price discovery can therefore be informative without being instantaneous or costless. Figure~\ref{fig:ch9-price-discovery} summarizes
the process through which information and investor
assessments are reflected in observed market prices,
with market frictions affecting this process.


\begin{figure}[H]

\centering

\begin{tikzpicture}[
  >=Latex,
  box/.style={
    draw,
    rounded corners,
    align=center,
    text width=2.55cm,
    minimum height=0.9cm
  },
  flow/.style={
    -Latex,
    line width=0.85pt,
    FinanceBlue
  }
]

\node[box] (info) at (0,1.4)
  {Public and private\\information};

\node[box] (beliefs) at (3.3,1.4)
  {Cash-flow and risk\\assessments};

\node[box] (orders) at (6.6,1.4)
  {Orders and\\trading};

\node[box] (price) at (9.9,1.4)
  {Observed\\market price};

\draw[flow] (info) -- (beliefs);
\draw[flow] (beliefs) -- (orders);
\draw[flow] (orders) -- (price);

\node[
  draw,
  rounded corners,
  align=center,
  minimum width=5.8cm,
  minimum height=0.75cm
] (frictions) at (6.6,-0.3)
  {Liquidity, funding, and market frictions};

\draw[flow] (frictions.north) -- (orders.south);

\end{tikzpicture}

\caption{Information, trading, and price discovery}
\label{fig:ch9-price-discovery}
\end{figure}

The efficient-market hypothesis provides a benchmark
for evaluating how information is incorporated into
market prices. Its implications depend on the
information set considered. Under the hypothesis,
investors should not be able to generate abnormal
risk-adjusted returns systematically using the
information covered by the relevant form of market
efficiency, after accounting for transaction and
information costs.

The hypothesis is commonly expressed in three forms according to the information set considered. Weak-form efficiency concerns information contained in past prices and trading activity. Semi-strong-form efficiency extends the information set to publicly available information, including financial reports and corporate announcements. Strong-form efficiency considers both public and private information and is best viewed as a demanding benchmark rather than a literal description of most markets.

Price discovery remains subject to economic frictions. Information can be costly to acquire, investors can disagree, and liquidity conditions, funding constraints, and limits to arbitrage can affect observed prices. If prices fully revealed all costly-to-acquire
information, investors would have little incentive
to incur the costs of acquiring it.

\begin{tcolorbox}[
  enhanced,
  breakable,
  colback=white,
  colframe=FinanceBlue,
  boxrule=0.7pt,
  sharp corners,
  borderline west={2.2pt}{0pt}{FinanceBlue},
  title=\textbf{Financial Perspective: Price Discovery Is Not Perfect Valuation},
  coltitle=white,
  colbacktitle=FinanceBlue,
  fonttitle=\bfseries,
  before skip=1.2em,
  after skip=1.2em
]

Observed market prices reflect investors' information,
expectations, and assessments under prevailing market
conditions. They provide signals about economic value,
but are not direct observations of an objectively
known fundamental value.

\end{tcolorbox}

\section{Equity Value and Risk}

Equity value reflects the present value of the cash flows available to shareholders, adjusted for the risk of those cash flows. Three elements are central to equity valuation: expected shareholder cash flows, the effect of reinvestment on future growth, and a required return consistent with equity risk.

\subsection{Equity Cash Flows and Valuation Inputs}

A distinction should be made between book equity and
market equity. Book equity is an accounting measure
derived from recognized assets and liabilities,
whereas market equity is the market value of
outstanding equity claims. The two need not coincide,
since market value reflects expectations about future
cash flows, growth opportunities, risk, and the rights
attached to the shares.

Dividends are direct cash payments from the corporation to shareholders. Share repurchases provide another form of payout: cash is paid to shareholders who sell their shares, while the number of shares outstanding declines. Dividends and repurchases therefore distribute value in different ways.

Retained earnings are not a distribution to shareholders when they are retained. Their contribution to equity value depends on the future cash flows generated by the resources kept within the firm. Retention creates value only to the extent that reinvestment supports sufficiently valuable future cash flows.

For an investor who plans to hold a share through date \(T\), the value today can be expressed as the present value of expected dividends during the holding period and the expected sale price:

\begin{equation}
P_0
=
\sum_{t=1}^{T}
\frac{\mathrm{E}[D_t]}{(1+r_e)^t}
+
\frac{\mathrm{E}[P_T]}{(1+r_e)^T},
\label{eq:ch9-finite-horizon-equity-value}
\end{equation}
where \(r_e\) is the required return on the equity claim under a constant-rate simplification.

If the terminal price is valued as the present value
of expected dividends beyond date $T$, extending
the valuation horizon indefinitely gives the
dividend discount representation:

\begin{equation}
P_0
=
\sum_{t=1}^{\infty}
\frac{\mathrm{E}[D_t]}{(1+r_e)^t}.
\label{eq:ch9-dividend-discount-model}
\end{equation}


For total equity valuation, analysts may instead forecast free cash flow to equity, \(FCFE_t\), the cash flow available to equity holders after operating needs, taxes, reinvestment, and net debt financing. A simplified valuation expression is

\begin{equation}
V_E
=
\sum_{t=1}^{\infty}
\frac{\mathrm{E}[FCFE_t]}{(1+r_e)^t}.
\label{eq:ch9-fcfe-valuation}
\end{equation}

The cash-flow measure and discount rate should refer to the same claim. Cash flows available to all capital providers require a discount rate appropriate to the entire enterprise, while cash flows available only to equity holders require an equity required return.

Market prices of comparable firms provide additional valuation benchmarks. Relative valuation uses measures such as price--earnings and price--book ratios to compare similar equity claims. Such comparisons are more informative when differences in growth, profitability, leverage, risk, and accounting measurement are taken into account.

\subsection{Growth Expectations and Market Value}

Growth affects equity value through its effect on future cash flows and the reinvestment required to generate them. Higher expected sales or earnings do not by themselves imply higher value. Reinvestment creates value when the return earned on additional equity exceeds the return required by shareholders.

The dividend payout ratio links expected earnings
to the next-period dividend. Suppose next-period
earnings per share are $E_1$, the firm retains a
fraction $k$ of earnings, and distributes the
remainder as dividends. Then

\begin{equation}
D_1=(1-k)E_1.
\label{eq:ch9-dividend-payout}
\end{equation}

If growth is financed through retained earnings, no
net new common equity is issued, and the return on
beginning book equity remains stable, a simplified
sustainable growth relation is

\begin{equation}
g=k\,ROE.
\label{eq:ch9-sustainable-growth}
\end{equation}

If dividends are expected to grow at a constant rate
$g$ indefinitely, the Gordon growth model, also known
as the constant-growth dividend discount model,
expresses the current share price as

\begin{equation}
P_0=\frac{D_1}{r_e-g},
\qquad
r_e>g.
\label{eq:ch9-gordon-growth}
\end{equation}

Substituting the payout and growth relations gives

\begin{equation}
P_0
=
\frac{(1-k)E_1}
{r_e-k\,ROE}.
\label{eq:ch9-growth-reinvestment-value}
\end{equation}

A higher retention ratio reduces current payout but can increase future growth. Its effect on value depends on the return earned on the retained resources relative to the required return on equity.

\begin{example}{Growth and Value Creation}{growth-reinvestment-value}

Suppose next-period earnings per share are 10, the retention ratio is 40\%, and the required return on equity is 10\%.

If retained earnings earn a stable \(ROE\) of 15\%,

\[
g=0.40\times 0.15=6\%,
\]
and
\[
P_0=\frac{(1-0.40)\times 10}{0.10-0.06}=\frac{6}{0.04}=150.
\]

If the return on equity is instead 8\%,

\[
g=0.40\times 0.08=3.2\%,
\]
and
\[
P_0=\frac{(1-0.40)\times 10}{0.10-0.032}=\frac{6}{0.068}\approx88.24.
\]

In this simplified setting, growth adds value when reinvestment earns more than the required return and reduces value when it earns less.

\end{example}

The same principle extends beyond constant-growth models. Rapid growth can require substantial investment, while modest current payout can be consistent with high value when retained resources are expected to generate sufficiently strong future cash flows. Growth expectations can also affect the timing of value. When a larger share of expected cash flows lies further in the future, equity value becomes more sensitive to changes in the required return.

\begin{tcolorbox}[
  enhanced,
  breakable,
  colback=white,
  colframe=FinanceBlue,
  boxrule=0.7pt,
  sharp corners,
  borderline west={2.2pt}{0pt}{FinanceBlue},
  title=\textbf{Financial Perspective: Growth Is Not the Same as Value Creation},
  coltitle=white,
  colbacktitle=FinanceBlue,
  fonttitle=\bfseries,
  before skip=1.2em,
  after skip=1.2em
]

Growth increases equity value only when the returns generated by the required reinvestment are sufficient relative to the return required by shareholders.

\end{tcolorbox}

\subsection{Equity Risk Premium and Required Return}

An investor's realized return on a share over one period combines the dividend received and the change in market price. If the initial price is \(P_0\), the end-of-period price is \(P_1\), and the dividend is \(D_1\), the realized holding-period return is

\begin{equation}
r
=
\frac{P_1-P_0+D_1}{P_0}.
\label{eq:ch9-equity-realized-return}
\end{equation}

A required return is different from a realized return. It is the return investors require for holding an equity claim given its risk and the returns available on alternative investments. Risks that can be diversified across many investments generally require less compensation than risks associated with broad economic conditions.

Let \(r_M\) denote the return on a broad market portfolio and \(r_f\) the risk-free rate. The market equity risk premium is

\begin{equation}
ERP
=
\mathrm{E}[r_M]-r_f.
\label{eq:ch9-equity-risk-premium}
\end{equation}

The equity risk premium is forward-looking. Historical averages of realized stock returns over risk-free returns can provide evidence about past outcomes, but they need not equal the premium investors require at a particular date.

The Capital Asset Pricing Model provides a widely used benchmark for relating a stock's required return to systematic risk. The beta of stock \(i\) is

\begin{equation}
\beta_i
=
\frac{\mathrm{Cov}(r_i,r_M)}
{\mathrm{Var}(r_M)}.
\label{eq:ch9-beta}
\end{equation}

Under the CAPM assumptions, the required return is

\begin{equation}
r_i
=
r_f
+
\beta_i\left(\mathrm{E}[r_M]-r_f\right).
\label{eq:ch9-capm}
\end{equation}

Beta measures the stock's exposure to market-wide return movements. A higher beta implies a larger required risk premium within the CAPM.

\begin{example}{A CAPM Required Return}{capm-required-return}

Suppose the risk-free rate is 3\%, the expected market equity risk premium is 6\%, and a stock has a beta of 1.2. The CAPM required return is

\[
r_i
=
3\%+1.2(6\%)
=
10.2\%.
\]

The estimate is conditional on the CAPM and on the inputs used.

\end{example}
The CAPM is a benchmark rather than a complete description of expected returns. Its broader economic insight is that risks linked to adverse market-wide conditions tend to require greater compensation because they are difficult to diversify when protection is most valuable. Beta represents this exposure through covariance with the market, while other asset-pricing models allow for additional sources of priced risk.

%========================================================
% Summary
%========================================================

\TailModule{summary}{Summary}

Common equity is a residual ownership claim. Shareholders receive value after senior contractual claims are satisfied and participate in the remaining upside, while limited liability ordinarily prevents losses on a fully paid share from extending beyond the value invested in that share. Equity financing provides capital without creating fixed contractual interest or principal payments and gives new investors residual cash-flow rights and, usually, voting rights. New issuance can dilute percentage ownership even when it increases the resources available to the firm.

Corporate governance allocates control over corporate resources. Shareholders usually exercise control indirectly through voting and the election of directors, while boards appoint and oversee management. Dispersed ownership can weaken monitoring incentives, whereas concentrated ownership can strengthen monitoring while increasing potential conflicts between controlling and minority shareholders. Share classes can separate cash-flow rights from voting rights, and corporate actions such as dividends, repurchases, stock splits, and new issues affect cash, share counts, and ownership claims in different ways.

Stock markets perform both financing and trading functions. Primary issues can bring new capital into the company, while secondary trades transfer existing shares among investors. Exchanges and other trading venues provide mechanisms through which orders interact, while liquidity depends on spreads, depth, price impact, information conditions, and the willingness of market participants to supply liquidity. Price discovery reflects information, expectations, risk assessments, and liquidity conditions in observable market prices through trading, without implying that every observed price equals an objectively known fundamental value.

Equity valuation requires consistency between shareholder cash flows and the required return used to discount them. Dividends and repurchases distribute cash, while retained earnings contribute to value through the future cash flows generated by reinvestment. Growth increases value when reinvestment is expected to earn sufficiently high returns relative to the required return on equity. The market equity risk premium represents the expected excess return on broad equity exposure, while the CAPM provides a benchmark linking an individual stock's required return to its systematic risk.

%========================================================
% Key Terms
%========================================================

\TailModule{keyterms}{Key Terms}

Common equity; residual claim; limited liability; ownership dilution; voting rights; board of directors; corporate governance; agency problem; controlling shareholder; minority shareholder; share class; preferred share; weighted voting rights; cash dividend; share repurchase; stock split; rights issue; initial public offering; seasoned equity offering; primary shares; secondary shares; exchange; alternative trading venue; market order; limit order; bid--ask spread; market liquidity; price discovery; informational efficiency; dividend discount model; free cash flow to equity; retention ratio; sustainable growth; equity risk premium; beta; Capital asset pricing model.

%========================================================
% Financial AI Lab
%========================================================

\TailModule{aiprompt}{Financial AI Lab}

\textbf{Lab: Audit an AI Equity Valuation and Governance Comparison}

Choose two publicly listed companies in the same broad industry. Ask an AI system to compare their equity valuations using a simple price multiple and either a free-cash-flow-to-equity or dividend-based approach. Require the system to identify the share class being valued, the relevant number of shares outstanding, the expected growth assumptions, the required return, and the source of each factual input.

Audit the result in four stages.

\begin{enumerate}

\item Verify whether the market capitalization and per-share valuation use the appropriate share count and whether each valuation multiple is matched to the corresponding equity measure.

\item Examine how the valuation responds to assumptions about growth and the required return. Test at least one alternative growth rate and one alternative required return or, when applicable, equity risk premium.

\item Compare the voting and cash-flow rights attached to the selected share classes. Determine whether one-share, one-vote applies or whether the governance structure separates voting power from economic ownership.

\item Check each factual claim against company disclosures or an official exchange source. Identify at least one statement for which the AI provides a plausible interpretation but insufficient supporting evidence.

\end{enumerate}

%========================================================
% Discussion Questions
%========================================================

\TailModule{questions}{Discussion Questions}

\begin{enumerate}

\item Why does limited liability not make common equity a low-risk claim? Distinguish between the maximum loss arising from ownership of the share itself and variation in the share's market value.

\item How can concentrated ownership reduce one agency problem while increasing another?

\item Why can a stock price rise sharply after favorable news even when no new cash enters the company from the stock market?

\item Why does faster earnings growth not necessarily imply a higher equity value? Relate your answer to reinvestment and the required return on equity.

\item Why should a historical average of realized market returns in excess of risk-free returns not be treated automatically as the current expected equity risk premium?

\end{enumerate}

%========================================================
% Applied Problems
%========================================================

\TailModule{appprobs}{Applied Problems}

\begin{enumerate}

\item \textbf{Residual payoffs and limited liability.}

A firm financed with debt and common equity will generate cash flow of either 160 or 70 next period, with each state equally likely. Promised debt repayment is 100.

\begin{enumerate}

\item Calculate the debt and equity payoffs in each state.

\item Calculate the expected payoff to equity.

\item Explain which part of the equity payoff structure reflects limited liability.

\item Explain why shareholders may benefit from a shift toward a riskier distribution of firm cash flows after debt has been issued, while creditors may lose.

\end{enumerate}

\item \textbf{Voting power and cash-flow ownership.}

A company has 20 million Class A shares with 10 votes
per share and 180 million Class B shares with one vote
per share. Assume that the two classes have identical
cash-flow rights per share.

\begin{enumerate}

\item Calculate the Class A share of the firm's total
cash-flow rights.

\item Calculate the Class A share of total voting power.

\item Explain why the two percentages differ.

\item Identify one potential benefit and one potential
governance risk of this structure.

\end{enumerate}

\item \textbf{Market and limit orders.}

The best bid for a stock is 24.96 for 6,000 shares and the best ask is 25.04 for 4,000 shares.

\begin{enumerate}

\item Calculate the quoted spread and midpoint.

\item Calculate the proportional quoted spread.

\item At what price can a market order to buy 2,000 shares execute if the quotation remains available?

\item If a limit order to buy 2,000 shares at 25.00 is submitted and no seller is willing to sell at 25.00 or below, how is the order treated in the limit order book?

\end{enumerate}

\item \textbf{Growth and reinvestment.}

A company expects next-period earnings per share of 8 and has a required return on equity of 11\%. The firm retains 30\% of earnings. Assume that growth is financed through retained earnings, no net new common equity is issued, and \(ROE\) remains stable at 16\%.

\begin{enumerate}

\item Calculate the sustainable growth rate.

\item Calculate next-period dividends per share.

\item Using the constant-growth model, calculate the share value.

\item Recalculate the value if \(ROE\) falls to 9\% while the retention ratio remains unchanged.

\item Explain the direction of the valuation change.

\end{enumerate}

\item \textbf{Equity risk and required return.}

The risk-free rate is 2.5\% and the expected market equity risk premium is 5.5\%.

\begin{enumerate}

\item Calculate the CAPM required return for a stock with beta 0.8.

\item Calculate the required return for a stock with beta 1.4.

\item Explain why the higher-beta stock has the higher required return within the CAPM.

\item Give one reason why an analyst might not want to rely on the CAPM estimate alone.

\end{enumerate}

\end{enumerate}

%========================================================
% Data Application
%========================================================

\TailModule{dataapp}{Data Application}

\textbf{Information and Price Discovery in China's Equity Market}

Select one A-share company listed on the Shanghai Stock Exchange or Shenzhen Stock Exchange and identify a dated corporate disclosure containing economically relevant new information, such as an earnings announcement, profit warning, major investment decision, dividend proposal, share repurchase plan, or new equity issue. Use the exchange disclosure platform and the company's formal announcement as the primary sources for the disclosure.

Collect the stock's daily closing prices for several trading days before and after the disclosure, together with a broad market index over the same period. Calculate daily stock returns and a simple market-adjusted abnormal return

\begin{equation*}
r_t=r_{i,t}-r_{M,t},
\end{equation*}

where \(r_{i,t}\) is the stock return and \(r_{M,t}\) is the market-index return on day \(t\).

Identify the exact time at which the information first became publicly available. If the disclosure was released after the market closed, treat the next trading day as the first day on which investors could respond through trading. Plot the stock price and report the market-adjusted returns around the disclosure. Determine whether the largest price movement occurred before, on, or after the first trading opportunity following the disclosure.

Interpret the result cautiously. A price reaction can reflect revisions to expected cash flows, required returns, or both. Market-wide news and other firm-specific information released near the same time can also affect observed returns. A short event window and daily closing prices do not by themselves establish causality or the speed of intraday price discovery.

Report the company, stock code, disclosure title, disclosure date and time, sample period, benchmark index, and data sources. Explain what additional evidence would be needed before concluding that the market anticipated the information, reacted slowly, or overreacted.

%========================================================
% Case Analysis
%========================================================

\TailModule{caseanalysis}{Case Analysis}

\CaseLabel{Historical Case Based on Public Information}

\textbf{Case: Xiaomi and Weighted Voting Rights in Hong Kong}

Xiaomi Corporation began trading in Hong Kong on July 9, 2018 and became the first company to list there with a weighted voting rights structure after the 2018 listing-regime reforms. Xiaomi's Class A shares carry ten votes per share on most shareholder resolutions, while Class B shares carry one vote per share. For specified reserved matters, each share carries one vote. The publicly listed shares are Class B shares.\footnote{Xiaomi Corporation, \emph{Global Offering Prospectus}, June 2018, \url{https://www1.hkexnews.hk/listedco/listconews/sehk/2018/0625/ltn20180625033.pdf}.}

The structure separates voting power from economic ownership. It can allow founders to retain substantial influence over corporate decisions after raising equity from public investors. Hong Kong's listing rules permit specified weighted-voting structures subject to eligibility requirements, continuing obligations, and safeguards for non-weighted-voting shareholders.

Analyze the structure from both financing and governance perspectives. Why might founders value the ability to raise public equity without surrendering proportional voting control? How could concentrated voting power support strategic continuity or long-horizon investment? What agency problems can arise when voting control exceeds cash-flow ownership? Which investor-protection mechanisms become more important when outside shareholders have weaker voting power?

Then consider the market-value implications. Explain why the market price of Xiaomi's publicly traded Class B shares reflects not only expected residual cash flows but also the governance rights and control structure associated with the investment. Would a change in voting rights necessarily leave market value unchanged if expected operating cash flows were initially unaffected? Explain how changes in control, agency costs, future cash-flow allocation, or the required return could affect equity value.

\emph{Primary sources: Hong Kong Exchanges and Clearing Limited, ``HKEX to Offer Investors More Options When Xiaomi Shares Debut,'' July 4, 2018; Xiaomi Corporation, 2018 Prospectus and 2018 Interim Report; Hong Kong Stock Exchange Listing Rules, Chapter 8A on Weighted Voting Rights.}

%========================================================
% Further Reading
%========================================================

\TailModule{furtherread}{Further Reading}

\begin{itemize}

\item Gordon, Myron J. (1959),
``Dividends, Earnings, and Stock Prices,''
\emph{The Review of Economics and Statistics}, 41(2, Part 1), 99--105.

\item Fama, Eugene F. (1970),
``Efficient Capital Markets: A Review of Theory and Empirical Work,''
\emph{The Journal of Finance}, 25(2), 383--417.

\item Grossman, Sanford J., and Joseph E. Stiglitz (1980),
``On the Impossibility of Informationally Efficient Markets,''
\emph{American Economic Review}, 70(3), 393--408.

\item Kyle, Albert S. (1985),
``Continuous Auctions and Insider Trading,''
\emph{Econometrica}, 53(6), 1315--1335.

\item La Porta, Rafael, Florencio Lopez-de-Silanes, Andrei Shleifer, and Robert W. Vishny (2000),
``Investor Protection and Corporate Governance,''
\emph{Journal of Financial Economics}, 58(1--2), 3--27.

\end{itemize}

\ChapterTail{summary,keyterms,aiprompt,questions,appprobs,dataapp,caseanalysis,furtherread}
